\section{Conclusions}
\label{sec:conclusion}

We presented \tool, an installable, extensible toolkit that makes adversarial evaluation of
graph-based spear-phishing and lateral-movement detectors reproducible and honest. By
coupling controllable cascade attacks with ground-truth labels and a leak-aware protocol
---benign traffic matched by design, victim-disjoint splits, a time-shuffle causality
control, and operational Recall@low-FPR metrics---\tool turns a landscape of incomparable,
confounder-prone results into a shared, falsifiable benchmark. Its single-interface design
lowers the cost of contributing new graphs, attacks, and detectors, its automatic red/blue
arms-race maps robustness frontiers instead of single points, and its fully synthetic
population makes the instrument freely distributable where real corpora cannot be.

The broader message is that, for graph-based targeted-attack detection, the bottleneck is
evaluation, not model capacity: node degree and send-time are exactly the shortcuts a
powerful model will exploit, so credible progress requires confounder controls and
operational metrics as first-class citizens. Future releases will broaden the graph library,
add further temporal GNN backbones, incorporate additional real communication graphs, and
package the arms-race search as a standing, versioned leaderboard so that the community can
accumulate comparable results over time.

\section*{Acknowledgements}
The authors thank colleagues at Kozminski University and NASK for feedback on the evaluation
design.

\section*{Code availability}
\tool is released under the MIT License at \url{https://github.com/kwarpechowski/cascadebench};
an archived snapshot with a citable Zenodo DOI will be deposited upon acceptance.